\chapter*{常见缩写和术语先导}
\addcontentsline{toc}{chapter}{常见缩写和术语先导}

第二册把程序放进现代硬件、并发运行时和分布式系统里解释。缩写和工程词很多，但它们都服务同一件事：把“慢、卡、错、丢、重复”变成可观察事实。

\begin{description}
  \item[CPU/GPU/OS/ISA/API/ABI] CPU 是中央处理器；GPU 是图形处理器或通用并行加速器；OS 是操作系统；ISA 是指令集架构；API 是源码层接口；ABI 是二进制接口。第二册会把同一段语义落到不同机器和调用边界上解释。
  \item[AMD] Advanced Micro Devices，常见 CPU/GPU 厂商名。本书提到 AMD 时，多数是在说明平台、微架构或工具事件不能和 Intel/Arm 混为一谈。
  \item[C++20/MSVC] C++20 是本书并发和同步示例使用的标准版本；MSVC 是 Microsoft Visual C++ 工具链。不同标准库和工具链的实现细节会影响源码阅读和实验结果。
  \item[CFG/STL/HPC] CFG 是控制流图；STL 是 C++ 标准库常见统称；HPC 是高性能计算。第二册会把这些词放进机器证据和工程实现里。
  \item[ASCII/UTF-8] ASCII 是早期单字节字符集合；UTF-8 是 Unicode 的常见变长字节编码。第二册解析器和 SIMD 字节分类会把“字节合同”和“文本语义”分开。
  \item[SIMD/ILP] SIMD 是一条指令处理多个数据；ILP 是 Instruction-Level Parallelism，指令级并行，让互不依赖的指令重叠执行。
  \item[AVX/AVX2/AVX-512] x86 常见向量扩展。AVX/AVX2 常见宽度是 256 位；AVX-512 提供更宽寄存器和 mask 机制，但要验证平台覆盖、频率和短输入成本。
  \item[VNNI/AMX/FMA] VNNI 更适合低精度点积；AMX 提供 tile 级矩阵能力；FMA 是 fused multiply-add，一条指令完成乘加。
  \item[ARM/POWER/SSE/LTO/PGO/JIT] ARM 和 POWER 是常见处理器架构生态；SSE 是较早的 x86 SIMD 扩展；LTO 是链接期优化；PGO 是 Profile-Guided Optimization，基于运行画像优化；JIT 是运行时即时编译或生成代码。
  \item[LLVM/GCC/SLP] LLVM 和 GCC 是常见编译器生态；SLP 是 Superword-Level Parallelism，编译器把多条相似标量语句合成向量操作的一类向量化方式。
  \item[perf/objdump/strace/numactl] 第二册常用实验工具。perf 采样和读取 PMU 计数器；objdump 看机器码和反汇编；strace 观察系统调用；numactl 控制或观察 NUMA 节点绑定。
  \item[PC/ALU/AGU/ROB/SIB/LEA] PC 是 program counter，程序计数器；ALU 是算术逻辑单元；AGU 是地址生成单元；ROB 是重排序缓冲区；SIB 是 x86 地址编码字段；LEA 是地址计算指令。它们分别解释取指位置、算术、地址和乱序窗口资源。
  \item[MMU/TLB/ITLB/DTLB/PTE/LLC] MMU 是内存管理单元；TLB 是地址翻译缓存；ITLB 缓存指令地址翻译，DTLB 缓存数据地址翻译；PTE 是页表项；LLC 是最后一级缓存。它们都是解释访存成本的核心证据来源。
  \item[L1/L2/L3/L1D] 缓存层级名。L1 通常最小最快，L2 居中，L3 常作为 LLC；L1D 特指一级数据缓存。
  \item[BTB/RAS] BTB 是 Branch Target Buffer，分支目标缓存；RAS 是 Return Address Stack，返回地址栈。它们帮助 CPU 猜下一条指令来自哪里。
  \item[CPI/p50/p95/p99] CPI 是每条指令平均周期数；p50、p95、p99 是分位数指标。并发性能不能只看平均值，p95/p99 往往暴露排队、迁核、缺页、锁竞争和 I/O 抖动。
  \item[MESI/MESIF/MOESI/HITM/LL/SC] MESI、MESIF、MOESI 是缓存一致性状态模型变体；HITM 表示命中被其他核心修改过的缓存行；LL/SC 是 load-linked/store-conditional 原子更新风格。
  \item[SRAM/RFO/THP/ULP] SRAM 是常用于 cache 的静态随机访问存储器；RFO 是 Read For Ownership，写入前获取缓存行所有权；THP 是透明大页；ULP 是浮点末位单位。
  \item[DRAM/VM/VMA/VFS/ASLR] DRAM 是主存；VM 是 virtual memory，虚拟内存；VMA 是进程虚拟内存区域；VFS 是 Linux 虚拟文件系统层；ASLR 是地址空间布局随机化。它们分别解释内存硬件、地址空间和文件系统抽象。
  \item[NUMA] Non-Uniform Memory Access，非统一内存访问。线程和内存不在同一节点时，远端访问会增加延迟和带宽压力。
  \item[PMU/TMA/IPC/CSV/GFLOP] PMU 是硬件性能计数器来源；TMA 是 Top-down Microarchitecture Analysis，一类把瓶颈分到前端、后端、错误推测和退休的分析方法；IPC 是每周期指令数；CSV 是实验表格文本；GFLOP/s 是每秒十亿次浮点运算。
  \item[SMT/COW/RMW] SMT 是同时多线程；COW 是写时复制；RMW 是读改写原子序列。并发性能经常卡在共享资源、共享页面或共享缓存行上。
  \item[TSO/TSAN/RAII] TSO 是 x86 常见的 Total Store Order 内存模型；TSAN 是 ThreadSanitizer，数据竞争检测工具；RAII 是 C++ 用对象生命周期管理资源的惯用法。
  \item[ASan/UBSan/sanitizer] ASan 查越界和 use-after-free；UBSan 查未定义行为；sanitizer 泛指这类插桩诊断工具。它们服务正确性，不代表发布性能。
  \item[CV/CFS/IRIW] CV 是 condition variable，条件变量；CFS 是 Linux Completely Fair Scheduler；IRIW 是 Independent Reads of Independent Writes，内存模型经典 litmus test。
  \item[ISR] Interrupt Service Routine，中断服务例程。它属于内核/硬件事件处理路径，实时系统和 I/O 讨论中会出现。
  \item[CAS] Compare-And-Swap，比较并交换。它是很多无锁结构的基础原子操作。
  \item[ABA] 无锁算法中的经典问题：一个位置从 A 变到 B 又变回 A，比较值相同但对象历史已经改变。
  \item[SPSC/MPSC/SPMC/MPMC] SPSC 是单生产者单消费者；MPSC 是多生产者单消费者；SPMC 是单生产者多消费者；MPMC 是多生产者多消费者。队列边界不同，原子变量和进度证明也不同。
  \item[EOF/EAGAIN/EINTR/ETIMEDOUT] EOF 表示输入结束；EAGAIN 表示稍后重试；EINTR 表示系统调用被信号中断；ETIMEDOUT 表示等待超时。并发和 I/O 状态机必须分别处理它们。
  \item[POSIX/epoll/SIGBUS] POSIX 是 Unix-like 系统接口规范；epoll 是 Linux 常见 I/O 事件通知机制；SIGBUS 是总线错误信号，mmap 文件被截断或映射生命周期错误时常见。
  \item[EFD/EPOLLIN/EPOLLOUT] EFD 常指出现在 eventfd flag 中的前缀；EPOLLIN 表示可读，EPOLLOUT 表示可写。它们都属于 Linux 事件驱动 I/O 的状态机输入。
  \item[EPOLLONESHOT/EPOLLRDHUP] epoll 事件标志。EPOLLONESHOT 表示一次事件后需要重新 armed；EPOLLRDHUP 表示对端半关闭，连接状态机必须处理。
  \item[SQE/CQE/SQ/CQ] io\_uring 的队列项和队列。SQE 是 submission queue entry，提交请求；CQE 是 completion queue entry，完成结果；SQ/CQ 是提交队列和完成队列。
  \item[EPOLLET/EPOLLERR/EPOLLHUP/SO\_ERROR] EPOLLET 是边沿触发模式；EPOLLERR 表示错误事件；EPOLLHUP 表示挂起或关闭；SO\_ERROR 用来读取 socket 上保存的异步错误。非阻塞 I/O 代码不能只处理“可读/可写”两个分支。
  \item[io\_uring] Linux 异步 I/O 接口。用户态把 SQE 放入提交队列，内核完成后写 CQE；要分清“提交成功”“I/O 完成”和“业务提交”三件事。
  \item[C++ 并发对象] \code{shared\_ptr} 是引用计数共享所有权；\code{lock\_guard} 是作用域锁；\code{co\_await} 是协程等待点。第二册看到它们时，要把语法还原成并发事实：谁拥有对象、谁持锁、谁可能暂停、谁负责恢复。
  \item[C 文件对象] \code{FILE} 是 C 标准库文件流类型。异步 I/O 章节出现 FILE 时，通常是为了和 file descriptor、socket、mmap 或 io\_uring 的完成语义对比；它不是文件名。
  \item[FIFO/LIFO/PID/TID] FIFO 是先进先出；LIFO 是后进先出；PID 是进程编号；TID 是线程编号。运行时和系统章节会用它们定位调度、线程和 I/O 对象。
  \item[RCU] Read-Copy-Update。读路径尽量无锁，更新路径复制新版本并延迟回收旧版本。
  \item[SLAB/SLUB] Linux 内核小对象分配器族，用对象缓存降低频繁分配释放成本。SLUB 是现代 Linux 常见实现。
  \item[TLS] Thread-Local Storage，线程局部存储。它能减少共享，但也会引入初始化、析构和汇总成本。
  \item[DMA/RDMA/UPI] DMA 让设备直接访问内存；RDMA 让远端机器的网卡直接读写内存区域；UPI 是 Intel 多 socket 互连之一，用于连接 CPU socket。
  \item[RPC] Remote Procedure Call，远程过程调用。看起来像本地函数调用，本质上跨机器、跨网络和失败边界。
  \item[RTT] Round-Trip Time，往返时间。分布式章节用它估算请求、复制和超时成本。
  \item[TCP/WAL/NFS] TCP 是可靠字节流传输协议；WAL 是 Write-Ahead Log，预写日志；NFS 是网络文件系统。分布式章节用它们解释网络边界、存储语义和崩溃恢复边界。
  \item[DAG] Directed Acyclic Graph，有向无环图。任务运行时常用 DAG 表达依赖关系。
  \item[BFS] Breadth-First Search，广度优先搜索。第二册提到 BFS 时，多用于图计算 kernel 的 frontier 扩展，而不是面试模板。
  \item[CSR/GEMM/GEMV] CSR 是压缩稀疏行格式；GEMM 是矩阵乘矩阵；GEMV 是矩阵乘向量。它们代表三类典型数据移动和计算形态。
  \item[CAP] Consistency、Availability、Partition tolerance 的分布式系统权衡框架。它不是设计清单，而是提醒网络分区下不能同时保证所有语义。
  \item[Raft/Paxos/quorum] Raft 和 Paxos 是共识协议族；quorum 是多数或法定票数集合。第二册用它们解释控制日志、leader、commit 和恢复边界。
  \item[leader term/fencing] leader term 是领导者任期；fencing 是用 term、epoch 或 lease version 阻止旧持有者继续写外部事实。
  \item[SLO] Service Level Objective，服务等级目标，例如 p95 延迟必须低于某个阈值。
  \item[P95/P99/RSS/GC/HDFS] P95/P99 是 95/99 分位延迟；RSS 是常驻内存大小；GC 是垃圾回收；HDFS 是 Hadoop Distributed File System。它们常出现在服务和分布式系统指标里。
  \item[GB/GFLOP/FLOP] GB 是十亿字节量级；FLOP 是浮点运算次数；GFLOP/s 是每秒十亿次浮点运算。Roofline 报告必须同时写字节和 FLOP 口径。
  \item[DWORD] double word，常见于 x86/Intel 风格汇编文本，通常表示 32 位内存操作数。
  \item[KV] Key/Value。第二册偶尔用 KV state 或 KV store 表示键值状态；第三册会把 KV cache 作为 Transformer 推理核心结构展开。
  \item[OOM] Out Of Memory，内存不足。并发和运行时章节遇到 OOM 时，要追问谁拥有内存、谁释放、失败能否恢复。
  \item[MR/NR/LSM/BLIS/XNNPACK/LMAX] MR/NR 是 GEMM microkernel 常见寄存器 tile 维度名；LSM 是日志结构合并树；BLIS 是高性能 BLAS 风格线性代数库；XNNPACK 是移动/边缘推理常见 CPU kernel 库；LMAX Disruptor 是序号环队列设计案例。本书提到它们时，是把工业项目当作源码阅读样本。
  \item[oneTBB] Intel oneAPI Threading Building Blocks，一套 C++ 并行运行时/算法库。第二册提到它时，重点是把任务、work stealing、arena 和调度证据和自写运行时对照。
  \item[CFS/EEVDF] Linux 调度器概念。CFS 是 Completely Fair Scheduler；EEVDF 是 newer Linux 中用于选择任务的虚拟截止期思想之一。读调度章节时，先把它们理解成“内核如何决定下一个线程运行”的策略。
  \item[FUTEX/FUTEX\_WAIT/FUTEX\_WAKE] futex 是 Linux 用户态锁和内核等待队列配合的基础接口。WAIT 负责睡眠等待某个内存值变化，WAKE 负责唤醒等待者；高层 mutex/condition variable 常会落到这类机制。
  \item[实验常量] \code{SUM_LANE_COUNT}、\code{CACHE_LINE_BYTES}、\code{EVENT_BUCKET_COUNT}、\code{BUFFER_SLOT_CAPACITY} 这类名字是实验或数据结构参数。它们把向量 lane 数、缓存行大小、事件桶数和环槽容量写清，便于复现实验。
  \item[CSV/JSON] CSV 是表格文本格式，适合保存实验数据；JSON 是结构化文本格式，适合保存配置、manifest 或报告。
  \item[worker/task/future/continuation] worker 是执行任务的线程或执行单元；task 是被调度的工作；future 是稍后可取得的结果；continuation 是前一步完成后继续执行的后续动作。
  \item[manifest/checkpoint/attempt/epoch] manifest 是机器可读清单；checkpoint 是可恢复状态；attempt 是同一个逻辑任务的一次物理尝试；epoch 是版本或轮次。它们共同解决“重试以后到底提交了什么”的问题。
  \item[owner/lease/fencing/snapshot] owner 是当前拥有资源或写权限的一方；lease 是有期限的租约；fencing 用版本 token 拦住旧 owner；snapshot 是某一时刻的状态副本。
  \item[partition/shard/replica/quorum] partition 是把数据或任务切分；shard 是切分后的分片；replica 是副本；quorum 是足以提交事实的法定票数集合。
  \item[batch/chunk/bucket/frontier] batch 是成批执行的一组输入；chunk 是连续切片；bucket 是按范围或哈希归类的桶；frontier 是图遍历当前要扩展的一层边界。
  \item[ring/head/tail/mask] ring 是环形缓冲；head/tail 是读写边界；mask 常用于把递增序号映射到环大小内的位置。不要把这些词当变量名，要看它们维护的队列不变量。
  \item[trace/profile/benchmark/baseline] trace 是按时间记录事件；profile 是运行画像；benchmark 是受控性能测量；baseline 是后续优化要对比的基线。
  \item[kernel/backend/runtime/scheduler] kernel 是热路径核心实现；backend 是落到具体硬件或系统接口的实现层；runtime 是负责调度、内存和同步的运行时；scheduler 是决定谁先执行、在哪里执行的策略。
  \item[layout/padding/stride/tile/lane] layout 是内存排列；padding 是填充；stride 是地址步幅；tile 是块化计算单元；lane 是 SIMD 向量中的元素槽。
  \item[backpressure/latency/throughput/bandwidth] backpressure 是下游处理不过来时让上游减速或失败；latency 是延迟；throughput 是吞吐；bandwidth 是单位时间数据搬运量。
  \item[block/stage/frame/pipeline] block 是一块数据、任务或代码基本块，具体含义看上下文；stage 是流水线阶段；frame 是栈帧、消息帧或数据帧；pipeline 是多个阶段串成的处理链。
  \item[reference/oracle/checksum/hash] reference 是定义正确结果的清晰实现；oracle 是判定输出是否正确的比较器；checksum 是结果摘要；hash 是把对象映射成摘要或桶位置的函数结果。
  \item[generation/watermark/metadata/schema/config] generation 是版本编号；watermark 是“已经安全推进到哪里”的水位；metadata 是描述数据的数据；schema 是结构规则；config 是配置。
  \item[register/stack/frame/exception/syscall] register 是寄存器；stack 是调用栈或栈式数据结构；frame 是一次调用或一帧数据；exception 是异常；syscall 是用户态进入内核的系统调用。
  \item[affinity/arena/span/fallback/handler] affinity 是线程、内存或中断绑定到特定 CPU/节点的亲和性；arena 是一块预规划内存；span 是非拥有连续视图；fallback 是保底路径；handler 是处理函数。
  \item[frontend/diagnostic/descriptor/target/pass] frontend 是输入导入和校验层；diagnostic 是结构化错误事实；descriptor 是描述对象；target 是构建目标或优化目标；pass 是分析或改写步骤。
  \item[IOMMU/MLP] IOMMU 是设备 DMA 访问内存时的地址翻译和隔离单元；MLP 在硬件性能语境常指 Memory-Level Parallelism，内存级并行，不是第三册 Transformer 里的多层感知机。
  \item[stream/prefill/tensor/replication/prefetch] stream 是异步事件流或执行队列；prefill 是模型服务先处理输入上下文的阶段；tensor 是带形状和元素类型的数据视图；replication 是复制副本；prefetch 是提前取数据以隐藏延迟。
  \item[coherence/shape/dtype/heap] coherence 是多核缓存一致性；shape 是数据形状或张量维度；dtype 是元素类型；heap 是动态分配区域。第二册会把这些词和 cache line、allocator、runtime state 放在一起看。
  \item[contract/callback/lowering/Raft] contract 是输入输出和失败边界；callback 是事件完成后的回调；lowering 是把高层描述降到更接近执行后端的表示；Raft 是共识协议族，用 leader、term、log 和 commit index 组织复制状态。
  \item[protocol version/magic/max frame] 分布式协议常见字段。version 表示协议版本，magic 用固定字节快速识别帧格式，max frame 限制单条消息大小，避免旧协议误读或恶意输入拖垮解析器。
  \item[trace 代号] \code{T2/A1}、\code{T1/A0} 这类短写通常是教学 trace 中的 task/attempt 代号。它们不是协议名，只是为了在日志表里压缩展示“第几个任务、第几次尝试”。
  \item[roofline] Roofline 是用算术强度、内存带宽和计算峰值估计性能上界的模型。它不是跑分图装饰，而是帮助判断一个 kernel 更受计算限制还是带宽限制。
  \item[数据布局缩写] AoS 是 array of structs，把每条记录的字段放在一起；SoA 是 struct of arrays，把同一字段集中成数组；AoSoA 是二者折中，按小块把字段重排。它们影响 cache line、SIMD 载入和预取效果，不是代码风格偏好。
  \item[C++ 原子接口] \code{fetch\_add} 是原子读改写加法；\code{memory\_order\_relaxed} 只保证原子性，不建立跨线程顺序；\code{memory\_order\_acquire/release} 常用于建立“发布数据、随后读取数据”的同步边。读原子代码必须先画出谁写、谁读、靠哪个顺序保证可见。
  \item[线程等待接口] \code{condition\_variable} 是条件变量；\code{notify\_one/notify\_all} 唤醒等待者；\code{try\_push}、\code{not\_full}、\code{not\_empty} 这类名字通常来自队列合同，分别表示尝试入队、队列未满、队列非空。它们的关键是不变量和唤醒条件，而不是函数名本身。
  \item[运行时对象名] \code{TaskState}、\code{FutureState}、\code{MessageBus}、\code{BatchView}、\code{RunReport}、\code{QueueContract} 这类 CamelCase 名字是本书实验框架对象。TaskState 保存任务状态，FutureState 保存异步结果，MessageBus 传递消息，BatchView 描述一批输入视图，RunReport 保存一次运行报告，QueueContract 描述队列语义。
  \item[分布式字段名] \code{task\_id}、\code{worker\_id}、\code{request\_id}、\code{record\_id}、\code{leader\_term}、\code{commit\_index} 这类字段不是缩写谜语，而是身份和版本边界。id 定位对象，term 区分领导者任期，commit index 表示日志已经提交到哪里。
  \item[开源系统名] DuckDB 是嵌入式分析数据库；ClickHouse 是列式分析数据库；OpenMP 是用编译指令表达并行的编程模型。附录把它们作为源码阅读对象，不要求先掌握全部内部实现。
  \item[实验短码] \code{T0/T1/T2}、\code{V0/V1}、\code{B0/B1}、\code{S0/S1} 这类短码通常是 trace、版本、批次或阶段编号。它们服务局部表格和日志压缩，不是新协议；读的时候看同一小节给它们绑定了哪个任务、版本或缓冲槽。
\end{description}
